\documentclass[UTF8]{ctexart}
\usepackage{amsmath,amsthm,amsfonts,amssymb}
\usepackage{mathtools}
\usepackage{mathrsfs}
\usepackage{bm}
\usepackage{extarrows}
\usepackage[a4paper, margin=1in]{geometry}
\usepackage{float}
\usepackage{indentfirst}
\usepackage{anyfontsize}
\usepackage{booktabs,multirow,multicol}
\usepackage[shortlabels,inline]{enumitem}
\usepackage{appendix}
\usepackage[dvipsnames]{xcolor}
\usepackage{graphicx}
\graphicspath{
    {./figure/}{./figures/}{./image/}{./images/}{./graphic/}{./graphics/}{./picture/}{./pictures/}
}
\usepackage{subcaption}

\usepackage[ruled,linesnumbered,noline]{algorithm2e}

\usepackage{hyperref}
\hypersetup{
    colorlinks=true,linkcolor=,urlcolor=cyan
}

\renewcommand*{\proofname}{\normalfont\bfseries Proof}

\usepackage{thmtools}

%% define environments

\declaretheorem[style=plain, name=定理, numbered=yes, numberwithin=section]{theorem}
\declaretheorem[style=plain, name=定理, numbered=no]{theorem*}

\declaretheorem[style=plain, name=命题, numbered=yes, sibling=theorem]{proposition}
\declaretheorem[style=plain, name=命题, numbered=no]{proposition*}

\declaretheorem[style=plain, name=推论, numbered=yes, sibling=theorem]{corollary}
\declaretheorem[style=plain, name=推论, numbered=no]{corollary*}

\declaretheorem[style=plain, name=引理, numbered=yes, sibling=theorem]{lemma}
\declaretheorem[style=plain, name=引理, numbered=no]{lemma*}

\declaretheorem[style=plain, name=陈述, numbered=yes, sibling=theorem]{claim}
\declaretheorem[style=plain, name=陈述, numbered=no]{claim*}

\declaretheorem[style=definition, name=定义, numbered=yes, numberwithin=section]{definition}
\declaretheorem[style=definition, name=定义, numbered=no]{definition*}

\declaretheorem[style=definition, name=例, numbered=yes, numberwithin=section]{example}
\declaretheorem[style=definition, name=例, numbered=no]{example*}

\declaretheorem[style=definition, name=习题, numbered=yes, numberwithin=section]{problem}
\declaretheorem[style=definition, name=习题, numbered=no]{problem*}

\declaretheorem[style=remark, name=注记, numbered=yes, numberwithin=section]{remark}
\declaretheorem[style=remark, name=注记, numbered=no]{remark*}

\declaretheoremstyle[headfont=\color{orange!80}\bfseries, bodyfont=\normalfont, spaceabove=3pt, spacebelow=3pt]{notestyle}

\declaretheorem[style=notestyle, name=笔记, numbered=yes, numberwithin=section]{note}
\declaretheorem[style=notestyle, name=笔记, numbered=no]{note*}

\declaretheoremstyle[headfont=\bfseries, bodyfont=\normalfont, spaceabove=3pt, spacebelow=3pt, qed=\ensuremath{\square}]{solutionstyle}

\declaretheorem[style=solutionstyle, name=解答, numbered=yes, numberwithin=section]{solution}
\declaretheorem[style=solutionstyle, name=解答, numbered=no]{solution*}

\usepackage[most]{tcolorbox}

\newcommand{\newtcbenvironment}[2]{
    \tcolorboxenvironment{#1}{#2, enhanced, breakable, sharp corners, boxrule=1pt}
    \tcolorboxenvironment{#1*}{#2, enhanced, breakable, rounded corners, boxrule=1pt}
}

%% define styles

\newtcbenvironment{theorem}{colframe=RoyalPurple, colback=RoyalPurple!8}
\newtcbenvironment{proposition}{colframe=RoyalPurple, colback=RoyalPurple!8}
\newtcbenvironment{corollary}{colframe=NavyBlue, colback=SkyBlue!8}
\newtcbenvironment{lemma}{colframe=NavyBlue, colback=SkyBlue!8}
\newtcbenvironment{claim}{colframe=NavyBlue, colback=SkyBlue!8}

\newtcbenvironment{definition}{colframe=ForestGreen, colback=ForestGreen!5}
\newtcbenvironment{example}{colframe=RawSienna, colback=RawSienna!5}
\newtcbenvironment{problem}{colframe=WildStrawberry!30, colback=WildStrawberry!5}

%% cbox
\newtcolorbox{cbox}[1][]{%
    enhanced,
    breakable,
    sharp corners,
    leftrule=2pt, rightrule=0pt, toprule=0pt, bottomrule=0pt,
    colframe=SkyBlue,
    colback=SkyBlue!8,
    #1
}

\title{LaTeX 笔记与定理模板使用说明}
\author{你的名字}
\date{\today}

\begin{document}

\maketitle

\section{简介}

本模板基于 \verb|ctexart| 文档类构建，专门为需要大量记录定理、定义和推导过程的理科笔记或作业设计。模板深度融合了 \verb|thmtools| 和 \verb|tcolorbox| 宏包，提供了一套色彩丰富、层次分明且易于扩展的定理环境。该说明文档由AI生成。

\section{核心特色与内容展示}

模板将不同的数学与逻辑环境进行了色彩分类，通过带颜色的背景框和边框来增强文档的视觉层次感。下面为您展示各个预设环境的实际渲染效果：

\subsection{核心结论（紫色系列）}
包含 \verb|theorem| (定理) 和 \verb|proposition| (命题)，适合放置最重要的核心结论。

\begin{theorem}[欧拉公式]
在复分析中，欧拉公式建立了三角函数与复指数函数之间的基本关系：
$$ e^{ix} = \cos x + i\sin x $$
\end{theorem}

\subsection{推导与辅助（蓝色系列）}
包含 \verb|corollary| (推论)、\verb|lemma| (引理) 和 \verb|claim| (陈述)，适合作为辅助性证明或衍生结论。

\begin{lemma}[辅助引理]
这里是引理的具体内容展示。辅助性推导通常使用蓝色系环境。
\end{lemma}

\subsection{概念引入（绿色系列）}
包含 \verb|definition| (定义)，用于引入新概念或专业术语。

\begin{definition*}[素数]
大于1的自然数中，除了1和它本身以外不再有其他因数的自然数。注意，这是一个带星号 (\verb|*|) 的环境，因此没有编号，且边框为圆角。
\end{definition*}

\subsection{实例与练习（棕/粉系列）}
包含 \verb|example| (例) 和 \verb|problem| (习题)。

\begin{example}[具体案例]
这是一个带有实例说明的环境展示。
\end{example}

\begin{problem}[课后思考]
证明上述欧拉公式在 $x = \pi$ 时的特殊情况（即欧拉恒等式）。
\end{problem}

\subsection{特殊文本与引用块}
包含无边框轻量环境（注记、笔记、解答）以及左侧高亮引用块 \verb|cbox|。

\begin{note*}[注意事项]
这是一条没有编号、带有橙色加粗标题的笔记。
\end{note*}

\begin{solution}
因为 $2$ 的因数只有 $1$ 和 $2$，所以 $2$ 是一个素数。注意末尾自动生成的 QED ($\square$) 符号。
\end{solution}

\begin{cbox}
\textbf{强调引用：} 这是一个侧边栏强调块 \verb|cbox|，非常适合用来做特别强调、补充说明或引用长文本。如果跨页内容较多，该盒子会自动分页。
\end{cbox}

\section{使用教程}

\subsection{编译环境}
由于本模板包含了大量的中文字符与高级字体/盒子设置，\textbf{必须使用 XeLaTeX 引擎进行编译} 才能保证正常输出。

\subsection{带星号 (*) 与不带星号的区别}
模板在定义盒子外观时进行了一项十分实用的设计：
\begin{itemize}
    \item \textbf{不带星号环境}（如 \verb|\begin{theorem}|）：环境\textbf{会自动编号}，并且盒子的四个角为\textbf{直角 (sharp corners)}。
    \item \textbf{带星号环境}（如 \verb|\begin{theorem*}|）：环境\textbf{不会进行编号}，并且盒子的四个角为\textbf{圆角 (rounded corners)}。
\end{itemize}

\subsection{如何自定义修改}
如果您需要在以后的笔记中修改颜色或新增定理环境（例如增加一个“假设”环境），请在导言区进行编辑：
\begin{enumerate}
    \item 在 \verb|%% define environments| 处使用 \verb|\declaretheorem| 定义环境的中文名称和编号逻辑。
    \item 在 \verb|%% define styles| 处使用 \verb|\newtcbenvironment| 为环境绑定所需的边框颜色 (\verb|colframe|) 和背景颜色 (\verb|colback|)。
\end{enumerate}

\end{document}